% 第9章《不定积分》题目解答
% 仅对应 chapters/ch09.tex 中的思考题、练习题与参考题。

\section*{第9章 不定积分：题目解答}

\paragraph{说明}
以下所有不定积分公式均理解为在被积函数定义域的任一连通区间上成立；若定义域有多个连通分支，各分支上的积分常数可以不同。涉及根式、对数等处，同时默认满足相应的实变量定义条件。

\subsection*{9.1.2 思考题}

\begin{enumerate}
  \item “不定积分”与“原函数”这两个概念有什么区别？有什么联系？

  \textbf{解}\quad 若 $F'(x)=f(x)$，则称 $F$ 是 $f$ 的一个原函数。原函数是一个具体函数；而不定积分
  \[
    \int f(x)\,\dd x
  \]
  表示 $f$ 的全部原函数所组成的函数族。在同一连通区间上，任意两个原函数之差为常数，因此若 $F$ 是一个原函数，则
  \[
    \int f(x)\,\dd x=F(x)+C.
  \]
  所以二者的联系是：求不定积分就是求原函数全体；区别是：一个原函数是函数族中的一个成员，而不定积分是整个函数族。

  \item 在 $x=0$ 点不连续的函数
  \[
    f(x)=
    \begin{cases}
      2x\sin\dfrac1x-\cos\dfrac1x,&x\ne0,\\
      0,&x=0
    \end{cases}
  \]
  在 $(-\infty,+\infty)$ 上是否有原函数？

  \textbf{解}\quad 有。定义
  \[
    F(x)=
    \begin{cases}
      x^2\sin\dfrac1x,&x\ne0,\\
      0,&x=0.
    \end{cases}
  \]
  当 $x\ne0$ 时，
  \[
    F'(x)=2x\sin\frac1x+x^2\cos\frac1x\left(-\frac1{x^2}\right)
    =2x\sin\frac1x-\cos\frac1x=f(x).
  \]
  而在 $x=0$，
  \[
    F'(0)=\lim_{h\to0}\frac{h^2\sin(1/h)}{h}
    =\lim_{h\to0}h\sin\frac1h=0=f(0).
  \]
  因此 $F$ 在整个实轴上可导且 $F'=f$，所以 $f$ 虽在 $0$ 点不连续，仍在 $(-\infty,+\infty)$ 上有原函数。

  \item 在公式
  \[
    \int\operatorname{sgn}x\,\dd x=\abs{x}+C,
    \qquad
    \int\frac1x\,\dd x=\ln\abs{x}+C
  \]
  中为什么会出现绝对值号？

  \textbf{解}\quad 两个公式都应按定义域的连通分支理解。对 $x>0$，
  \[
    \operatorname{sgn}x=1,
    \qquad
    \frac1x=(\ln x)',
  \]
  因而原函数分别为 $x+C$ 与 $\ln x+C$；对 $x<0$，
  \[
    \operatorname{sgn}x=-1,
    \qquad
    \frac1x=(\ln(-x))',
  \]
  因而原函数分别为 $-x+C$ 与 $\ln(-x)+C$。把正、负两个分支统一写出，就得到 $\abs{x}$ 与 $\ln\abs{x}$。

  还要注意，$\operatorname{sgn}x$ 在 $0$ 处有第一类间断点，因此它在整个 $\RR$ 上没有原函数；公式 $\abs{x}+C$ 只是在 $(-\infty,0)$ 和 $(0,+\infty)$ 上分别成立。$1/x$ 本身也只在这两个区间上有定义。

  \item 下列等式是否正确？说明理由。
  \begin{enumerate}[label=(\arabic*)]
    \item $\dd\int f(x)\,\dd x=f(x)$；

    \textbf{解}\quad 不正确。左边是对原函数族中的每个函数取微分所得的微分式，而右边是函数。正确的形式应为
    \[
      \dd\int f(x)\,\dd x=f(x)\,\dd x.
    \]

    \item $\dd\int f(x)\,\dd x=f(x)\,\dd x$；

    \textbf{解}\quad 正确。若 $F'(x)=f(x)$，则
    \[
      \dd(F(x)+C)=F'(x)\,\dd x=f(x)\,\dd x.
    \]

    \item $\int\dd f(x)=f(x)$；

    \textbf{解}\quad 不正确。不定积分表示全部原函数，因此
    \[
      \int\dd f(x)=f(x)+C,
    \]
    不能漏掉任意常数。

    \item $\dd\int\dd f(x)=\dd f(x)$。

    \textbf{解}\quad 正确。因为
    \[
      \int\dd f(x)=f(x)+C,
    \]
    对其中每个函数再取微分都得到 $\dd f(x)$。
  \end{enumerate}

  \item 以下推导中有什么错误？
  \[
    \int\frac{\dd x}{x}
    =x\cdot\frac1x-\int x\,\dd\left(\frac1x\right)
    =1+\int\frac{\dd x}{x},
  \]
  因而推出 $0=1$。

  \textbf{解}\quad 错误在于把两个不定积分中的任意常数当成了同一个常数，从而非法相消。不定积分等式表示的是原函数族。写明常数后应为
  \[
    \ln\abs{x}+C_1
    =1+\ln\abs{x}+C_2,
  \]
  这要求 $C_1=C_2+1$，完全没有矛盾。也就是说，分部积分公式成立到一个任意常数；不能把等式两侧出现的同形不定积分当成同一个固定函数再直接消去。
\end{enumerate}

\subsection*{9.1.6 练习题}

\begin{enumerate}
  \item 计算下列不定积分。
  \begin{enumerate}[label=(\arabic*)]
    \item
    \[
      \int\frac{x}{1+x^4}\,\dd x.
    \]
    令 $u=x^2$，则
    \[
      \int\frac{x}{1+x^4}\,\dd x
      =\frac12\int\frac{\dd u}{1+u^2}
      =\frac12\arctan x^2+C.
    \]

    \item
    \[
      \int\frac{\dd x}{\sqrt{x(1-x)}}.
    \]
    在 $0<x<1$ 上令 $x=\sin^2t$，则
    \[
      \dd x=2\sin t\cos t\,\dd t,
      \qquad
      \sqrt{x(1-x)}=\sin t\cos t,
    \]
    故
    \[
      \int\frac{\dd x}{\sqrt{x(1-x)}}
      =2t+C
      =2\arcsin\sqrt{x}+C.
    \]

    \item
    \[
      \int\ln(1+x^2)\,\dd x.
    \]
    分部积分得
    \begin{align*}
      \int\ln(1+x^2)\,\dd x
      &=x\ln(1+x^2)-\int\frac{2x^2}{1+x^2}\,\dd x\\
      &=x\ln(1+x^2)-2x+2\arctan x+C.
    \end{align*}

    \item
    \[
      \int\frac{\arctan\sqrt{x}}{\sqrt{x}(1+x)}\,\dd x.
    \]
    对 $x>0$ 令 $t=\sqrt{x}$，则 $\dd x=2t\,\dd t$，所以
    \[
      \int\frac{\arctan\sqrt{x}}{\sqrt{x}(1+x)}\,\dd x
      =2\int\frac{\arctan t}{1+t^2}\,\dd t
      =(\arctan\sqrt{x})^2+C.
    \]

    \item
    \[
      \int\frac{\dd x}{\ee^x-1}.
    \]
    令 $u=\ee^x$，则
    \[
      \int\frac{\dd x}{\ee^x-1}
      =\int\frac{\dd u}{u(u-1)}
      =\ln\abs{u-1}-\ln u+C,
    \]
    因而
    \[
      \boxed{\int\frac{\dd x}{\ee^x-1}=\ln\abs{1-\ee^{-x}}+C}.
    \]

    \item
    \[
      \int\frac{x\ln(x+\sqrt{1+x^2})}{(1+x^2)^2}\,\dd x.
    \]
    记
    \[
      u=\ln(x+\sqrt{1+x^2}),
      \qquad
      \dd u=\frac{\dd x}{\sqrt{1+x^2}},
    \]
    且
    \[
      \frac{x\,\dd x}{(1+x^2)^2}
      =-\frac12\,\dd\left(\frac1{1+x^2}\right).
    \]
    分部积分得
    \begin{align*}
      \int\frac{x u}{(1+x^2)^2}\,\dd x
      &=-\frac{u}{2(1+x^2)}
        +\frac12\int\frac{\dd x}{(1+x^2)^{3/2}}\\
      &=-\frac{u}{2(1+x^2)}+\frac{x}{2\sqrt{1+x^2}}+C.
    \end{align*}
    即
    \[
      \boxed{-\frac{\ln(x+\sqrt{1+x^2})}{2(1+x^2)}
      +\frac{x}{2\sqrt{1+x^2}}+C}.
    \]

    \item
    \[
      \int\frac{x\ee^x}{(1+x)^2}\,\dd x.
    \]
    注意
    \[
      \left(\frac{\ee^x}{1+x}\right)'
      =\frac{x\ee^x}{(1+x)^2},
    \]
    所以
    \[
      \boxed{\int\frac{x\ee^x}{(1+x)^2}\,\dd x=\frac{\ee^x}{1+x}+C}.
    \]

    \item
    \[
      \int\frac{1+x}{x(1+x\ee^x)}\,\dd x.
    \]
    乘以 $\ee^x/\ee^x$，并令 $u=x\ee^x$，有
    \[
      \dd u=\ee^x(1+x)\,\dd x,
    \]
    因而
    \begin{align*}
      \int\frac{1+x}{x(1+x\ee^x)}\,\dd x
      &=\int\frac{\dd u}{u(1+u)}\\
      &=\ln\abs{u}-\ln\abs{1+u}+C\\
      &=x+\ln\abs{x}-\ln(1+x\ee^x)+C.
    \end{align*}

    \item
    \[
      \int\ln^2(x+\sqrt{1+x^2})\,\dd x.
    \]
    令
    \[
      u=\ln(x+\sqrt{1+x^2}),
      \qquad
      x=\sinh u,
      \qquad
      \dd x=\cosh u\,\dd u.
    \]
    则
    \begin{align*}
      \int u^2\cosh u\,\dd u
      &=u^2\sinh u-2\int u\sinh u\,\dd u\\
      &=u^2\sinh u-2u\cosh u+2\sinh u+C.
    \end{align*}
    因而
    \[
      \boxed{x\ln^2(x+\sqrt{1+x^2})
      -2\sqrt{1+x^2}\ln(x+\sqrt{1+x^2})+2x+C}.
    \]

    \item
    \[
      \int\frac{\ee^{\arctan x}}{(1+x^2)^{3/2}}\,\dd x.
    \]
    令 $t=\arctan x$，则 $x=\tan t$，$\dd x=(1+x^2)\,\dd t$，且
    \[
      \frac1{\sqrt{1+x^2}}=\cos t.
    \]
    因此
    \begin{align*}
      \int\frac{\ee^{\arctan x}}{(1+x^2)^{3/2}}\,\dd x
      &=\int\ee^t\cos t\,\dd t\\
      &=\frac12\ee^t(\sin t+\cos t)+C\\
      &=\boxed{\frac{(x+1)\ee^{\arctan x}}{2\sqrt{1+x^2}}+C}.
    \end{align*}
  \end{enumerate}

  \item 用配对积分法计算下列不定积分。
  \begin{enumerate}[label=(\arabic*)]
    \item
    \[
      \int\frac{\dd x}{1+x^2+x^4}.
    \]
    因
    \[
      1+x^2+x^4=(x^2+x+1)(x^2-x+1),
    \]
    且
    \[
      \frac1{1+x^2+x^4}
      =\frac12\frac{x+1}{x^2+x+1}
       +\frac12\frac{1-x}{x^2-x+1},
    \]
    积分得
    \[
      \boxed{
      \frac14\ln\frac{x^2+x+1}{x^2-x+1}
      +\frac1{2\sqrt3}\left[
        \arctan\frac{2x+1}{\sqrt3}
        +\arctan\frac{2x-1}{\sqrt3}
      \right]+C}.
    \]

    \item
    \[
      \int\frac{\ee^x\,\dd x}{\ee^x+\ee^{-x}}.
    \]
    乘以 $\ee^x/\ee^x$，得
    \[
      \int\frac{\ee^{2x}}{1+\ee^{2x}}\,\dd x
      =\boxed{\frac12\ln(1+\ee^{2x})+C}.
    \]

    \item
    \[
      \int\frac{b\sin x+a\cos x}{a\sin x+b\cos x}\,\dd x,
      \qquad a\ne b.
    \]
    设
    \[
      D=a\sin x+b\cos x,
      \qquad
      D'=a\cos x-b\sin x.
    \]
    因 $a\ne b$ 保证 $(a,b)\ne(0,0)$，有
    \[
      b\sin x+a\cos x
      =\frac{2ab}{a^2+b^2}D
       +\frac{a^2-b^2}{a^2+b^2}D'.
    \]
    因此
    \[
      \boxed{
      \frac{2ab}{a^2+b^2}x
      +\frac{a^2-b^2}{a^2+b^2}\ln\abs{a\sin x+b\cos x}+C}.
    \]

    \item
    \[
      \int\frac{\dd x}{1+x^3}.
    \]
    分解
    \[
      \frac1{1+x^3}
      =\frac1{3(x+1)}+\frac{-x+2}{3(x^2-x+1)},
    \]
    故
    \[
      \boxed{
      \frac13\ln\abs{x+1}
      -\frac16\ln(x^2-x+1)
      +\frac1{\sqrt3}\arctan\frac{2x-1}{\sqrt3}+C}.
    \]
  \end{enumerate}

  \item 通过两个配对积分求 $\int\cos^4x\,\dd x$ 与 $\int\sin^4x\,\dd x$。

  \textbf{解}\quad 先有
  \[
    \cos^4x-\sin^4x=(\cos^2x-\sin^2x)(\cos^2x+\sin^2x)=\cos2x,
  \]
  因而
  \[
    A:=\int(\cos^4x-\sin^4x)\,\dd x
    =\frac12\sin2x+C.
  \]
  又
  \[
    \cos^4x+\sin^4x
    =1-2\sin^2x\cos^2x
    =\frac34+\frac14\cos4x,
  \]
  所以
  \[
    B:=\int(\cos^4x+\sin^4x)\,\dd x
    =\frac34x+\frac1{16}\sin4x+C.
  \]
  取线性组合可得
  \[
    \boxed{
    \int\cos^4x\,\dd x
    =\frac38x+\frac14\sin2x+\frac1{32}\sin4x+C},
  \]
  \[
    \boxed{
    \int\sin^4x\,\dd x
    =\frac38x-\frac14\sin2x+\frac1{32}\sin4x+C}.
  \]

  \item 导出计算下列不定积分的递推公式。
  \begin{enumerate}[label=(\arabic*)]
    \item 令
    \[
      I_n=\int\sin^nx\,\dd x.
    \]
    对 $n>1$，分部积分
    \begin{align*}
      I_n
      &=-\int\sin^{n-1}x\,\dd(\cos x)\\
      &=-\sin^{n-1}x\cos x
        +(n-1)\int\sin^{n-2}x\cos^2x\,\dd x\\
      &=-\sin^{n-1}x\cos x+(n-1)I_{n-2}-(n-1)I_n.
    \end{align*}
    因此
    \[
      \boxed{I_n=-\frac1n\sin^{n-1}x\cos x+\frac{n-1}{n}I_{n-2}}.
    \]
    可取 $I_0=x$、$I_1=-\cos x$ 为初始项。

    \item 令
    \[
      T_n=\int\tan^nx\,\dd x.
    \]
    对 $n>1$，由 $\tan^2x=\sec^2x-1$，
    \begin{align*}
      T_n
      &=\int\tan^{n-2}x(\sec^2x-1)\,\dd x\\
      &=\frac{\tan^{n-1}x}{n-1}-T_{n-2}.
    \end{align*}
    即
    \[
      \boxed{T_n=\frac{\tan^{n-1}x}{n-1}-T_{n-2}}.
    \]

    \item 令
    \[
      S_n=\int\sec^nx\,\dd x.
    \]
    对 $n>1$，分部积分
    \begin{align*}
      S_n
      &=\int\sec^{n-2}x\,\dd(\tan x)\\
      &=\sec^{n-2}x\tan x
        -(n-2)\int\sec^{n-2}x\tan^2x\,\dd x\\
      &=\sec^{n-2}x\tan x-(n-2)(S_n-S_{n-2}),
    \end{align*}
    因而
    \[
      \boxed{S_n=\frac{\sec^{n-2}x\tan x}{n-1}
      +\frac{n-2}{n-1}S_{n-2}}.
    \]

    \item 令
    \[
      J_n=\int\frac{\dd x}{x^n\sqrt{1+x^2}}.
    \]
    计算
    \begin{align*}
      \frac{\dd}{\dd x}\left(\frac{\sqrt{1+x^2}}{x^{n-1}}\right)
      &=(2-n)\frac1{x^{n-2}\sqrt{1+x^2}}
       +(1-n)\frac1{x^n\sqrt{1+x^2}}.
    \end{align*}
    对 $n>1$ 积分整理得
    \[
      \boxed{J_n
      =-\frac{\sqrt{1+x^2}}{(n-1)x^{n-1}}
       -\frac{n-2}{n-1}J_{n-2}}.
    \]
    可用
    \[
      J_0=\ln(x+\sqrt{1+x^2})+C,
      \qquad
      J_1=\ln\abs{\frac{x}{1+\sqrt{1+x^2}}}+C
    \]
    作为初始积分。
  \end{enumerate}

  \item 试求。
  \begin{enumerate}[label=(\arabic*)]
    \item
    \[
      \int x f''(x)\,\dd x.
    \]
    分部积分得
    \[
      \boxed{\int x f''(x)\,\dd x=x f'(x)-f(x)+C}.
    \]

    \item
    \[
      \int f'(2x)\,\dd x.
    \]
    令 $u=2x$，则
    \[
      \boxed{\int f'(2x)\,\dd x=\frac12f(2x)+C}.
    \]
  \end{enumerate}

  \item 设
  \[
    I(m,n)=\int\cos^mx\sin^nx\,\dd x.
  \]
  证明下列递推关系。
  \begin{enumerate}[label=(\arabic*)]
    \item 证明
    \[
      I(m,n)=-\frac{\sin^{n-1}x\cos^{m+1}x}{m+n}
      +\frac{n-1}{m+n}I(m,n-2).
    \]
    从
    \[
      I(m,n)=-\int\cos^mx\sin^{n-1}x\,\dd(\cos x)
    \]
    出发，分部积分得
    \begin{align*}
      I(m,n)
      &=-\sin^{n-1}x\cos^{m+1}x
       +\int\cos x\,\dd(\cos^mx\sin^{n-1}x)\\
      &=-\sin^{n-1}x\cos^{m+1}x
       -mI(m,n)
       +(n-1)\int\cos^{m+2}x\sin^{n-2}x\,\dd x.
    \end{align*}
    又
    \[
      \int\cos^{m+2}x\sin^{n-2}x\,\dd x
      =I(m,n-2)-I(m,n),
    \]
    因而
    \[
      (m+n)I(m,n)
      =-\sin^{n-1}x\cos^{m+1}x+(n-1)I(m,n-2),
    \]
    即得所证公式。

    \item 证明
    \[
      I(n,n)=-\frac{\cos2x\sin^{n-1}2x}{n\,2^{n+1}}
      +\frac{n-1}{4n}I(n-2,n-2).
    \]
    由
    \[
      \sin x\cos x=\frac12\sin2x
    \]
    得
    \[
      I(n,n)=\frac1{2^n}\int\sin^n2x\,\dd x
      =\frac1{2^{n+1}}\int\sin^nu\,\dd u,
      \qquad u=2x.
    \]
    使用本题第 4(1) 小题的递推式，
    \begin{align*}
      I(n,n)
      &=-\frac{\cos2x\sin^{n-1}2x}{n\,2^{n+1}}
       +\frac{n-1}{n\,2^{n+1}}\int\sin^{n-2}u\,\dd u.
    \end{align*}
    而
    \[
      I(n-2,n-2)=\frac1{2^{n-1}}\int\sin^{n-2}u\,\dd u,
    \]
    故
    \[
      \frac{n-1}{n\,2^{n+1}}\int\sin^{n-2}u\,\dd u
      =\frac{n-1}{4n}I(n-2,n-2),
    \]
    结论成立。
  \end{enumerate}
\end{enumerate}

\subsection*{9.2.4 练习题}

\begin{enumerate}
  \item 用观察法将被积函数拆开后计算不定积分。
  \begin{enumerate}[label=(\arabic*)]
    \item
    \[
      \frac{x}{(x+1)(x+2)}=-\frac1{x+1}+\frac2{x+2},
    \]
    所以
    \[
      \boxed{-\ln\abs{x+1}+2\ln\abs{x+2}+C}.
    \]

    \item
    \[
      \frac{x^2+x+1}{x(1+x^2)}=\frac1x+\frac1{1+x^2},
    \]
    因而
    \[
      \boxed{\ln\abs{x}+\arctan x+C}.
    \]

    \item
    \[
      \frac1{x^4(x^2+1)}=\frac1{x^4}-\frac1{x^2}+\frac1{x^2+1},
    \]
    故
    \[
      \boxed{-\frac1{3x^3}+\frac1x+\arctan x+C}.
    \]

    \item
    \[
      \frac{4x^2+3}{(x^2+1)(x^2+2)}
      =-\frac1{x^2+1}+\frac5{x^2+2},
    \]
    因而
    \[
      \boxed{-\arctan x+\frac5{\sqrt2}\arctan\frac{x}{\sqrt2}+C}.
    \]
  \end{enumerate}

  \item 计算下列有理函数的不定积分。
  \begin{enumerate}[label=(\arabic*)]
    \item
    \[
      \int\frac{x^5-x}{1+x^8}\,\dd x.
    \]
    令 $t=x^2$，则
    \[
      \int\frac{x^5-x}{1+x^8}\,\dd x
      =\frac12\int\frac{t^2-1}{1+t^4}\,\dd t.
    \]
    利用
    \[
      \int\frac{t^2-1}{1+t^4}\,\dd t
      =\frac1{2\sqrt2}
       \ln\frac{t^2-\sqrt2t+1}{t^2+\sqrt2t+1}+C,
    \]
    得
    \[
      \boxed{
      \frac1{4\sqrt2}
      \ln\frac{x^4-\sqrt2x^2+1}{x^4+\sqrt2x^2+1}+C}.
    \]

    \item
    \[
      \int\frac{\dd x}{x(x^n+a)},\qquad a\ne0.
    \]
    分解为
    \[
      \frac1{x(x^n+a)}
      =\frac1a\left(\frac1x-\frac{x^{n-1}}{x^n+a}\right),
    \]
    所以
    \[
      \boxed{
      \frac1a\ln\abs{x}-\frac1{an}\ln\abs{x^n+a}+C}.
    \]

    \item
    \[
      \int\frac{x\,\dd x}{(x+1)(x^2+3)}.
    \]
    有
    \[
      \frac{x}{(x+1)(x^2+3)}
      =-\frac1{4(x+1)}+\frac{x+3}{4(x^2+3)},
    \]
    因而
    \[
      \boxed{-\frac14\ln\abs{x+1}
      +\frac18\ln(x^2+3)
      +\frac{\sqrt3}{4}\arctan\frac{x}{\sqrt3}+C}.
    \]

    \item
    \[
      \int\frac{x-1}{(x^2+2x+3)^2}\,\dd x.
    \]
    令 $u=x+1$，则分母为 $(u^2+2)^2$，分子为 $u-2$。利用
    \[
      \int\frac{\dd u}{(u^2+2)^2}
      =\frac{u}{4(u^2+2)}+\frac1{4\sqrt2}\arctan\frac{u}{\sqrt2},
    \]
    可得
    \[
      \boxed{-\frac{x+2}{2(x^2+2x+3)}
      -\frac1{2\sqrt2}\arctan\frac{x+1}{\sqrt2}+C}.
    \]

    \item
    \[
      \int\frac{1+x+x^2}{(x-2)^{10}}\,\dd x.
    \]
    令 $u=x-2$，则 $x^2+x+1=u^2+5u+7$，故
    \begin{align*}
      \int\frac{1+x+x^2}{(x-2)^{10}}\,\dd x
      &=\int(u^{-8}+5u^{-9}+7u^{-10})\,\dd u\\
      &=\boxed{-\frac1{7u^7}-\frac5{8u^8}-\frac7{9u^9}+C},
    \end{align*}
    其中 $u=x-2$。

    \item
    \[
      \int\frac{x^3+1}{x^3-5x^2+6x}\,\dd x.
    \]
    先作除法并分解：
    \[
      \frac{x^3+1}{x(x-2)(x-3)}
      =1+\frac1{6x}-\frac9{2(x-2)}+\frac{28}{3(x-3)}.
    \]
    因此
    \[
      \boxed{x+\frac16\ln\abs{x}
      -\frac92\ln\abs{x-2}
      +\frac{28}{3}\ln\abs{x-3}+C}.
    \]

    \item
    \[
      \int\frac{\dd x}{1+x^6}.
    \]
    因
    \[
      1+x^6=(x^2+1)(x^2+\sqrt3x+1)(x^2-\sqrt3x+1),
    \]
    且
    \[
      \frac1{1+x^6}
      =\frac1{3(x^2+1)}
       +\frac{\sqrt3x+2}{6(x^2+\sqrt3x+1)}
       -\frac{\sqrt3x-2}{6(x^2-\sqrt3x+1)},
    \]
    所以
    \[
      \boxed{
      \frac{\sqrt3}{12}\ln\frac{x^2+\sqrt3x+1}{x^2-\sqrt3x+1}
      +\frac16\left[
        \arctan(2x+\sqrt3)+\arctan(2x-\sqrt3)
      \right]
      +\frac13\arctan x+C}.
    \]

    \item
    \[
      \int\frac{x^2-x+3}{(x^2+x+1)(x-1)^2}\,\dd x.
    \]
    分解为
    \[
      \frac{x^2-x+3}{(x^2+x+1)(x-1)^2}
      =\frac1{(x-1)^2}-\frac2{3(x-1)}
       +\frac{2(x+2)}{3(x^2+x+1)}.
    \]
    积分得
    \[
      \boxed{-\frac1{x-1}-\frac23\ln\abs{x-1}
      +\frac13\ln(x^2+x+1)
      +\frac2{\sqrt3}\arctan\frac{2x+1}{\sqrt3}+C}.
    \]
  \end{enumerate}

  \item 计算下列三角函数有理式的不定积分。
  \begin{enumerate}[label=(\arabic*)]
    \item
    \[
      \int\frac{\dd x}{1+\cos x}
      =\frac12\int\sec^2\frac x2\,\dd x
      =\boxed{\tan\frac x2+C}.
    \]

    \item
    \[
      \int\frac{\dd x}{\sin x\cos^4x}.
    \]
    令 $u=\cos x$，则 $\dd u=-\sin x\,\dd x$，从而
    \[
      \int\frac{\dd x}{\sin x\cos^4x}
      =-\int\frac{\dd u}{u^4(1-u^2)}.
    \]
    又
    \[
      \frac1{u^4(1-u^2)}=\frac1{u^4}+\frac1{u^2}+\frac1{1-u^2},
    \]
    因此
    \[
      \boxed{
      \frac13\sec^3x+\sec x+\ln\abs{\tan\frac x2}+C}.
    \]

    \item
    \[
      \int\frac{\dd x}{2+\tan^2x}.
    \]
    令 $t=\tan x$，则
    \[
      \frac{\dd x}{2+\tan^2x}
      =\frac{\dd t}{(1+t^2)(2+t^2)}
      =\left(\frac1{1+t^2}-\frac1{2+t^2}\right)\dd t.
    \]
    故在任一不跨越 $\tan x$ 极点的区间上
    \[
      \boxed{x-\frac1{\sqrt2}\arctan\frac{\tan x}{\sqrt2}+C}.
    \]

    \item
    \[
      \int\tan^5x\sec^3x\,\dd x.
    \]
    令 $u=\sec x$，$\dd u=\sec x\tan x\,\dd x$，则
    \begin{align*}
      \int\tan^5x\sec^3x\,\dd x
      &=\int\tan^4x\sec^2x\,\dd u\\
      &=\int u^2(u^2-1)^2\,\dd u.
    \end{align*}
    因而
    \[
      \boxed{\frac17\sec^7x-\frac25\sec^5x+\frac13\sec^3x+C}.
    \]

    \item
    \[
      \int\frac{\sec x}{(1+\sec x)^2}\,\dd x.
    \]
    令 $t=\tan(x/2)$，利用
    \[
      \cos x=\frac{1-t^2}{1+t^2},
      \qquad
      \dd x=\frac{2\,\dd t}{1+t^2},
    \]
    得
    \[
      \frac{\sec x}{(1+\sec x)^2}\,\dd x
      =\frac12(1-t^2)\,\dd t.
    \]
    所以
    \[
      \boxed{\frac12\tan\frac x2-\frac16\tan^3\frac x2+C}.
    \]

    \item
    \[
      \int\frac{\sin^2x\cos x}{\sin x+\cos x}\,\dd x.
    \]
    一个便于检验的分解结果为
    \[
      \frac{\sin^2x\cos x}{\sin x+\cos x}
      =\frac{\dd}{\dd x}\left[
        \frac14\ln\abs{\sin x+\cos x}
        -\frac14\cos x(\sin x+\cos x)
      \right].
    \]
    因而
    \[
      \boxed{\frac14\ln\abs{\sin x+\cos x}
      -\frac14\cos x(\sin x+\cos x)+C}.
    \]

    \item
    \[
      \int\frac{\sin x\cos x}{1+\sin^4x}\,\dd x.
    \]
    令 $u=\sin^2x$，则 $\dd u=2\sin x\cos x\,\dd x$，故
    \[
      \boxed{\frac12\arctan(\sin^2x)+C}.
    \]

    \item
    \[
      \int\frac{\dd x}{\sin^4x\cos^2x}.
    \]
    令 $t=\cot x$，则 $\dd x=-\dd t/(1+t^2)$，并有
    \[
      \frac{\dd x}{\sin^4x\cos^2x}
      =-\left(t^{-2}+2+t^2\right)\dd t.
    \]
    因此
    \[
      \boxed{\tan x-2\cot x-\frac13\cot^3x+C}.
    \]
  \end{enumerate}

  \item 计算 Poisson（泊松）积分
  \[
    \int\frac{1-r^2}{1-2r\cos x+r^2}\,\dd x,
    \qquad -1<r<1.
  \]
  令 $t=\tan(x/2)$。由于
  \[
    \cos x=\frac{1-t^2}{1+t^2},
    \qquad
    \dd x=\frac{2\,\dd t}{1+t^2},
  \]
  分母化为
  \[
    1-2r\cos x+r^2
    =\frac{(1-r)^2+(1+r)^2t^2}{1+t^2}.
  \]
  因此
  \begin{align*}
    \int\frac{1-r^2}{1-2r\cos x+r^2}\,\dd x
    &=2(1-r^2)\int
      \frac{\dd t}{(1-r)^2+(1+r)^2t^2}\\
    &=\boxed{2\arctan\left(
      \frac{1+r}{1-r}\tan\frac x2
      \right)+C}.
  \end{align*}
  该表达式在每个不跨越 $x=(2k+1)\pi$ 的区间上直接成立；跨越这些点时只需调整反正切的分支常数即可得到全局原函数。

  \item 计算下列无理函数的不定积分。
  \begin{enumerate}[label=(\arabic*)]
    \item
    \[
      \int\frac{\dd x}{\sqrt{\sin x\cos^7x}}.
    \]
    在原根式为实数且非零的区间上有 $\tan x>0$。令
    \[
      t=\sqrt{\tan x},
      \qquad \tan x=t^2.
    \]
    则
    \[
      \dd x=\frac{2t}{1+t^4}\,\dd t,
      \qquad
      \sqrt{\sin x\cos^7x}=\frac{t}{(1+t^4)^2},
    \]
    从而
    \[
      \int\frac{\dd x}{\sqrt{\sin x\cos^7x}}
      =2\int(1+t^4)\,\dd t.
    \]
    因此
    \[
      \boxed{2\sqrt{\tan x}+\frac25\tan^{5/2}x+C}.
    \]

    \item
    \[
      \int\frac{\dd x}{x\sqrt{2x^2+3}}.
    \]
    令 $t=\sqrt{2x^2+3}$。由 $t^2-3=2x^2$ 得
    \[
      \frac{\dd x}{x t}=\frac{\dd t}{t^2-3},
    \]
    所以
    \[
      \boxed{
      \frac1{2\sqrt3}\ln\abs{
        \frac{\sqrt{2x^2+3}-\sqrt3}{\sqrt{2x^2+3}+\sqrt3}}
      +C}.
    \]

    \item
    \[
      \int\sqrt{\tan^2x+2}\,\dd x.
    \]
    先令 $t=\tan x$，得到
    \[
      \int\frac{\sqrt{t^2+2}}{1+t^2}\,\dd t.
    \]
    再令 $t=\sqrt2\sinh u$，则
    \[
      \frac{\sqrt{t^2+2}}{1+t^2}\,\dd t
      =\left(1+\frac1{\cosh2u}\right)\dd u.
    \]
    利用
    \[
      \int\frac{\dd u}{\cosh2u}
      =\frac12\arctan(\sinh2u)+C,
    \]
    且
    \[
      u=\ln\frac{t+\sqrt{t^2+2}}{\sqrt2},
      \qquad
      \sinh2u=t\sqrt{t^2+2},
    \]
    得
    \[
      \boxed{
      \ln\frac{\tan x+\sqrt{\tan^2x+2}}{\sqrt2}
      +\frac12\arctan\left(\tan x\sqrt{\tan^2x+2}\right)+C}.
    \]

    \item
    \[
      \int\frac{x\ee^x\,\dd x}{\sqrt{1+\ee^x}}.
    \]
    因
    \[
      \frac{\ee^x\,\dd x}{\sqrt{1+\ee^x}}
      =2\,\dd\sqrt{1+\ee^x},
    \]
    分部积分得
    \[
      I=2x\sqrt{1+\ee^x}-2\int\sqrt{1+\ee^x}\,\dd x.
    \]
    再令 $u=\sqrt{1+\ee^x}$，则
    \[
      \int\sqrt{1+\ee^x}\,\dd x
      =\int\frac{2u^2}{u^2-1}\,\dd u
      =2u+\ln\abs{\frac{u-1}{u+1}}+C.
    \]
    因而
    \[
      \boxed{
      2x\sqrt{1+\ee^x}-4\sqrt{1+\ee^x}
      -2\ln\abs{\frac{\sqrt{1+\ee^x}-1}{\sqrt{1+\ee^x}+1}}+C}.
    \]

    \item
    \[
      \int\frac{\dd x}{\sqrt{(x-1)^3(x-2)}}.
    \]
    实根式要求 $x<1$ 或 $x>2$。在任一这样的区间上令
    \[
      t=\sqrt{\frac{x-2}{x-1}}.
    \]
    由
    \[
      x-1=\frac1{1-t^2},
      \qquad
      \dd x=\frac{2t}{(1-t^2)^2}\,\dd t,
    \]
    可得被积式恰化为 $2\,\dd t$，故
    \[
      \boxed{2\sqrt{\frac{x-2}{x-1}}+C}.
    \]

    \item
    \[
      \int\frac{\sqrt{x}-1}{\sqrt[3]{x}+1}\,\dd x.
    \]
    对 $x>0$ 令 $t=x^{1/6}$，则 $x=t^6$、$\dd x=6t^5\,\dd t$，于是
    \[
      I=6\int\frac{t^5(t^3-1)}{t^2+1}\,\dd t.
    \]
    多项式除法给出
    \[
      \frac{t^5(t^3-1)}{t^2+1}
      =t^6-t^4-t^3+t^2+t-1-\frac{t-1}{t^2+1}.
    \]
    所以
    \[
      \boxed{
      6\left(
        \frac{t^7}{7}-\frac{t^5}{5}-\frac{t^4}{4}
        +\frac{t^3}{3}+\frac{t^2}{2}-t
        -\frac12\ln(1+t^2)+\arctan t
      \right)+C,
      \quad t=x^{1/6}}.
    \]
    该原函数可由连续性延拓到 $x=0$。
  \end{enumerate}
\end{enumerate}

\subsection*{9.3.2 参考题}

\begin{enumerate}
  \item 设
  \[
    f'(\sin^2x)=\cos4x+\tan^2x,
    \qquad 0<x<1,
  \]
  求 $f$。

  \textbf{解}\quad 令 $y=\sin^2x$。由于 $0<x<1$，故 $0<y<\sin^21$。并且
  \[
    \cos4x=1-8y(1-y)=1-8y+8y^2,
    \qquad
    \tan^2x=\frac{y}{1-y}=-1+\frac1{1-y}.
  \]
  因而
  \[
    f'(y)=8y^2-8y+\frac1{1-y}.
  \]
  对 $y$ 积分，得
  \[
    f(y)=\frac83y^3-4y^2-\ln(1-y)+C.
  \]
  所以在 $0<y<\sin^21$ 上
  \[
    \boxed{f(y)=\frac83y^3-4y^2-\ln(1-y)+C}.
  \]
  若把自变量仍记为 $x$，则为
  \[
    \boxed{f(x)=\frac83x^3-4x^2-\ln(1-x)+C}.
  \]

  \item 计算下列不定积分。
  \begin{enumerate}[label=(\arabic*)]
    \item
    \[
      I=\int x\arctan x\ln(1+x^2)\,\dd x.
    \]
    取
    \[
      u=\arctan x,
      \qquad
      \dd v=x\ln(1+x^2)\,\dd x.
    \]
    有
    \[
      v=\frac12\left[(1+x^2)\ln(1+x^2)-(1+x^2)\right].
    \]
    因此
    \begin{align*}
      I
      &=\frac12(1+x^2)[\ln(1+x^2)-1]\arctan x\\
      &\quad-\frac12\int[\ln(1+x^2)-1]\,\dd x.
    \end{align*}
    又
    \[
      \int\ln(1+x^2)\,\dd x
      =x\ln(1+x^2)-2x+2\arctan x+C,
    \]
    故
    \[
      \boxed{
      \frac12(1+x^2)[\ln(1+x^2)-1]\arctan x
      -\frac x2\ln(1+x^2)+\frac32x-\arctan x+C}.
    \]

    \item
    \[
      \int\frac{1-\ln x}{(x-\ln x)^2}\,\dd x.
    \]
    对 $x>0$，直接计算
    \[
      \left(\frac{x}{x-\ln x}\right)'
      =\frac{1-\ln x}{(x-\ln x)^2},
    \]
    故
    \[
      \boxed{\frac{x}{x-\ln x}+C}.
    \]

    \item
    \[
      \int [x]\abs{\sin\pi x}\,\dd x,
      \qquad x\ge0.
    \]
    令
    \[
      n=\lfloor x\rfloor,
      \qquad
      u=x-n\in[0,1).
    \]
    在 $[n,n+1)$ 上有 $[x]=n$ 且
    \[
      \abs{\sin\pi x}=\sin\pi u.
    \]
    从 $0$ 积到 $x$，前面每个完整区间 $[k,k+1]$ 的贡献为 $2k/\pi$，所以
    \begin{align*}
      F(x)
      &=\sum_{k=0}^{n-1}\frac{2k}{\pi}
        +\frac{n}{\pi}(1-\cos\pi u)\\
      &=\frac{n^2-n\cos\pi u}{\pi}.
    \end{align*}
    因而一个在 $[0,+\infty)$ 上连续可导的原函数为
    \[
      \boxed{
      F(x)=\frac{n^2-n\cos\pi(x-n)}{\pi}+C,
      \qquad n=\lfloor x\rfloor}.
    \]

    \item
    \[
      \int\sin x\ln(\sin x)\,\dd x.
    \]
    实对数要求 $\sin x>0$。令 $u=\cos x$，则
    \[
      I=-\frac12\int\ln(1-u^2)\,\dd u.
    \]
    分部积分可得
    \[
      \int\ln(1-u^2)\,\dd u
      =u\ln(1-u^2)-2u+\ln\abs{\frac{1+u}{1-u}}+C.
    \]
    代回 $u=\cos x$，并用 $\ln(1-\cos^2x)=2\ln(\sin x)$，得到
    \[
      \boxed{
      \cos x\,[1-\ln(\sin x)]
      +\ln\abs{\tan\frac x2}+C}.
    \]

    \item
    \[
      \int\frac{\dd x}{\sin(x+a)\sin(x+b)}.
    \]
    若 $\sin(b-a)\ne0$，利用
    \[
      \cot(x+a)-\cot(x+b)
      =\frac{\sin(b-a)}{\sin(x+a)\sin(x+b)},
    \]
    得
    \[
      \boxed{
      \frac1{\sin(b-a)}
      \ln\abs{\frac{\sin(x+a)}{\sin(x+b)}}+C}.
    \]
    若 $b-a=k\pi$，则
    \[
      \sin(x+b)=(-1)^k\sin(x+a),
    \]
    因而
    \[
      \boxed{
      \int\frac{\dd x}{\sin(x+a)\sin(x+b)}
      =(-1)^{k+1}\cot(x+a)+C}.
    \]

    \item
    \[
      \int\frac{1-x^n}{x(1+x^n)}\,\dd x,
      \qquad n\in\NN_+.
    \]
    因
    \[
      \frac{1-x^n}{x(1+x^n)}
      =\frac1x-\frac{2x^{n-1}}{1+x^n},
    \]
    故
    \[
      \boxed{\ln\abs{x}-\frac2n\ln\abs{1+x^n}+C}.
    \]
  \end{enumerate}

  \item 对每个正整数 $n$，定义
  \[
    I_n=\int\frac{(ax+b)^n}{\sqrt{cx+b}}\,\dd x,
  \]
  求递推公式。

  \textbf{解}\quad 记
  \[
    y=ax+b,
    \qquad z=cx+b.
  \]
  对 $c\ne0$，计算
  \[
    \dd(y^n\sqrt z)
    =na y^{n-1}\sqrt z\,\dd x
     +\frac c2\frac{y^n}{\sqrt z}\,\dd x.
  \]
  因此
  \[
    cI_n=2y^n\sqrt z-2na\int y^{n-1}\sqrt z\,\dd x.
  \]
  又有恒等式
  \[
    az=cy+b(a-c),
  \]
  从而
  \[
    a\int y^{n-1}\sqrt z\,\dd x
    =cI_n+b(a-c)I_{n-1}.
  \]
  代入上式并整理，得到
  \[
    \boxed{
    I_n=\frac{2(ax+b)^n\sqrt{cx+b}}{c(2n+1)}
    -\frac{2n\,b(a-c)}{c(2n+1)}I_{n-1}},
    \qquad c\ne0.
  \]
  初始项为
  \[
    I_0=\frac2c\sqrt{cx+b}+C.
  \]
  若 $c=0$ 且根式有定义，则 $\sqrt{cx+b}=\sqrt b$ 为常数，积分退化为多项式积分，可直接计算，无须递推。

  \item 对于实数 $a\ne0$ 与正整数 $n>2$，定义
  \[
    I_n=\int\left[
      \frac{\sin((x-a)/2)}{\sin((x+a)/2)}
    \right]^n\,\dd x.
  \]
  证明题中递推公式。

  \textbf{证明}\quad 记
  \[
    y=\frac{\sin((x-a)/2)}{\sin((x+a)/2)}.
  \]
  在分母非零的区间上，直接求导得
  \[
    y'=\frac{\sin a}{2\sin^2((x+a)/2)}.
  \]
  另一方面，利用两角相差 $a$，可验证
  \[
    y^2-2\cos a\,y+1
    =\frac{\sin^2a}{\sin^2((x+a)/2)}.
  \]
  因此
  \begin{align*}
    y^n+y^{n-2}-2\cos a\,y^{n-1}
    &=y^{n-2}\frac{\sin^2a}{\sin^2((x+a)/2)}\\
    &=2\sin a\,y^{n-2}y'\\
    &=\frac{2\sin a}{n-1}(y^{n-1})'.
  \end{align*}
  两边积分即得
  \[
    I_n+I_{n-2}-2\cos a\,I_{n-1}
    =\frac{2\sin a}{n-1}y^{n-1},
  \]
  即
  \[
    \boxed{
    I_n=-I_{n-2}+2\cos a\,I_{n-1}
    +\frac{2\sin a}{n-1}
    \left[
      \frac{\sin((x-a)/2)}{\sin((x+a)/2)}
    \right]^{n-1}}.
  \]

  \item 设 $Q$ 为 $n$ 次多项式，具有 $n$ 个相异实根 $x_i$；$P$ 与 $Q$ 不可约，且 $\deg P=m<n$。证明
  \[
    \int\frac{P(x)}{Q(x)}\,\dd x
    =\sum_{i=1}^n\frac{P(x_i)}{Q'(x_i)}\ln\abs{x-x_i}+C.
  \]

  \textbf{证明}\quad 因 $Q$ 的根均为单根，可写成
  \[
    Q(x)=q_n\prod_{i=1}^n(x-x_i).
  \]
  由部分分式分解，存在唯一常数 $A_i$ 使
  \[
    \frac{P(x)}{Q(x)}=\sum_{i=1}^n\frac{A_i}{x-x_i}.
  \]
  对固定的 $i$，两边乘以 $x-x_i$ 并令 $x\to x_i$，得
  \[
    A_i=\lim_{x\to x_i}\frac{(x-x_i)P(x)}{Q(x)}.
  \]
  令 $Q(x)=(x-x_i)Q_i(x)$，则
  \[
    Q_i(x_i)=Q'(x_i),
  \]
  所以
  \[
    A_i=\frac{P(x_i)}{Q'(x_i)}.
  \]
  因而
  \[
    \frac{P(x)}{Q(x)}
    =\sum_{i=1}^n\frac{P(x_i)}{Q'(x_i)}\frac1{x-x_i}.
  \]
  在任一不含根 $x_i$ 的连通区间上逐项积分，即得
  \[
    \boxed{
    \int\frac{P(x)}{Q(x)}\,\dd x
    =\sum_{i=1}^n\frac{P(x_i)}{Q'(x_i)}\ln\abs{x-x_i}+C}.
  \]

  \item 求 $\int f(x)\,\dd x$，其中 $f(x)$ 为 $x$ 到离其最近的整数的距离。

  \textbf{解}\quad 令
  \[
    n=\lfloor x\rfloor,
    \qquad
    u=x-n\in[0,1).
  \]
  则
  \[
    f(x)=
    \begin{cases}
      u,&0\le u\le\dfrac12,\\[2mm]
      1-u,&\dfrac12\le u<1.
    \end{cases}
  \]
  该函数以 $1$ 为周期，且每个完整周期的积分为
  \[
    \int_0^1f(x)\,\dd x
    =2\int_0^{1/2}u\,\dd u=\frac14.
  \]
  因此可取原函数
  \[
    \boxed{
    F(x)=\frac n4+
    \begin{cases}
      \dfrac{u^2}{2},&0\le u\le\dfrac12,\\[2mm]
      u-\dfrac{u^2}{2}-\dfrac14,&\dfrac12\le u<1,
    \end{cases}
    \quad n=\lfloor x\rfloor,\ u=x-n}.
  \]
  在半整数和整数连接点处，左右函数值及导数均吻合，因此这是整个实轴上的原函数；最后再加任意常数 $C$。

  \item 使用题给 Liouville 定理证明
  \[
    \int\ee^{-x^2}\,\dd x,
    \qquad
    \int\frac{\ee^x}{x}\,\dd x
  \]
  均不是初等不定积分。

  \textbf{证明}\quad 先考虑
  \[
    \int\ee^{-x^2}\,\dd x.
  \]
  若它是初等函数，则由题给定理，存在有理函数 $h$ 使
  \[
    \int\ee^{-x^2}\,\dd x=h(x)\ee^{-x^2}+C.
  \]
  求导得到
  \[
    h'(x)-2xh(x)=1.
  \]
  假设有理函数 $h$ 有有限极点 $x=a$，阶数为 $m\ge1$。则 $h'$ 在 $a$ 处有 $m+1$ 阶极点，而 $2xh$ 至多有 $m$ 阶极点（当 $a=0$ 时甚至更低），因此上述等式左边必有极点，不可能等于常数 $1$。所以 $h$ 没有有限极点，只能是多项式。

  若 $h$ 是非零的 $d$ 次多项式，则 $-2xh$ 的次数为 $d+1$，而 $h'$ 的次数至多为 $d-1$，最高次项不可能抵消，故该等式不可能成立；$h=0$ 也不满足 (1)。矛盾。因此
  \[
    \boxed{\int\ee^{-x^2}\,\dd x\text{ 不是初等函数}}.
  \]

  再考虑
  \[
    \int\frac{\ee^x}{x}\,\dd x.
  \]
  若它是初等函数，则存在有理函数 $h$ 使
  \[
    \int\frac{\ee^x}{x}\,\dd x=h(x)\ee^x+C.
  \]
  求导得到
  \[
    h'(x)+h(x)=\frac1x.
  \]
  若 $h$ 在某个 $a\ne0$ 处有 $m$ 阶极点，则 $h'$ 有 $m+1$ 阶极点，不能被 $h$ 抵消，而等式右边在 $a$ 处解析，矛盾。因此 $h$ 的有限极点只能可能在 $0$。

  若 $h$ 在 $0$ 有 $m\ge1$ 阶极点，则 $h'$ 有 $m+1$ 阶极点，而 $h$ 只有 $m$ 阶极点，左边最高阶极点仍无法抵消；这也不可能等于只有一阶极点的 $1/x$。若 $h$ 在 $0$ 无极点，则左边在 $0$ 解析，同样不可能等于 $1/x$。故不存在这样的有理函数 $h$，于是
  \[
    \boxed{\int\frac{\ee^x}{x}\,\dd x\text{ 不是初等函数}}.
  \]

  最后，若 $\int\dfrac{\dd x}{\ln x}$ 是初等函数，作代换 $x=\ee^t$ 就会得到初等原函数
  \[
    \int\frac{\ee^t}{t}\,\dd t,
  \]
  与上面的结论矛盾。因此 $\int\dfrac{\dd x}{\ln x}$ 也不是初等不定积分。
\end{enumerate}
